\section{Introdução}
\label{sec:introducao}
O presente relatório descreve o desenvolvimento do primeiro projeto prático da disciplina de Computação Gráfica. O trabalho aborda um dos conceitos fundamentais da computação gráfica bidimensional: a transformação entre diferentes sistemas de coordenadas utilizados durante o processo de visualização.

Nesse processo, um ponto inicialmente definido no Sistema de Coordenadas do Mundo é transformado para o Sistema de Coordenadas Normalizadas do Dispositivo (NDC) e, posteriormente, para as Coordenadas do Dispositivo, que correspondem à posição do ponto na tela em pixels.

O projeto consiste no desenvolvimento de uma aplicação em Java capaz de realizar e demonstrar essas transformações matemáticas. A aplicação utiliza a biblioteca Java Swing para apresentar os resultados graficamente, permitindo visualizar o mapeamento de um ponto do espaço do mundo até sua representação na tela.
